\chapter{Delayed Sleep-Wake Phase Disorder}
\index{Delayed Sleep-Wake Phase Disorder}
\label{sec:Delayed Sleep-Wake Phase Disorder}

\section{Overview}
Prevalence is 1--3\% in adults and 7--16\% in adolescents, making it the most common circadian rhythm disorder. This condition is among the most common mental health disorders worldwide. It affects people across all age groups, socioeconomic backgrounds, and geographic regions.
\section{Causes}
\section{Risk Factors}

\begin{itemize}[noitemsep]
  \item Genetic predisposition and family history
  \item Advancing age
  \item Lifestyle factors (diet, exercise, smoking, alcohol)
  \item Environmental exposures
  \item Underlying medical conditions
  \item Medication use
\end{itemize}
\section{Signs and Symptoms}

\subsection{Common symptoms}
\begin{itemize}[noitemsep]
  \item \item You may experience severe difficulty falling asleep at conventional times but normal sleep once initiated when allowed to follow the endogenous schedule. Pain or discomfort in the affected area    \item Pain or discomfort in the affected area
  \end{itemize}

\subsection{Emergency warning signs}
\begin{itemize}[noitemsep]
  \item Severe or worsening symptoms of delayed sleep-wake phase disorder require immediate medical evaluation
\end{itemize}
\section{Progression and Stages}
Delayed sleep-wake phase disorder (DSWPD) is a circadian rhythm sleep-wake disorder characterized by a stable delay of the sleep-wake phase relative to the desired or required schedule, with sleep onset typically after 1:00--3:00 AM and wake times in the late morning or early afternoon. This condition happens because of an abnormally long intrinsic circadian period, reduced phase-advance capacity, or a hypo-responsive phase-advance mechanism to morning light. The condition may progress through different stages of severity over time. Early stages often involve mild or subtle symptoms that may be overlooked. Moderate stages involve more noticeable symptoms affecting daily function. Advanced stages may involve significant impairment and complications.
\section{Complications}
\begin{itemize}[noitemsep]
  \item Disease progression leading to worsening symptoms
  \item Secondary infections or injuries
  \item Side effects of treatment
  \item Reduced quality of life
  \item Emotional and psychological impact
\end{itemize}
\section{Diagnosis}
Doctors diagnose this using clinical history, sleep diaries, actigraphy, and dim-light melatonin onset measurement. Comprehensive medical history and physical examination Laboratory tests to assess organ function and rule out other conditions Imaging studies to visualize affected structures Specialized testing depending on the affected system Consultation with relevant specialists Monitoring of disease progression over time

\begin{itemize}[noitemsep]
  \item Comprehensive medical history and physical examination
  \item Laboratory tests to assess organ function
  \item Imaging studies to visualize affected structures
  \item Specialized testing depending on the affected system
  \item Consultation with relevant specialists
\end{itemize}
\section{Prevention}
\begin{itemize}[noitemsep]
  \item Maintain a healthy lifestyle with balanced diet and exercise
  \item Avoid known risk factors
  \item Attend regular medical check-ups for early detection
  \item Follow recommended screening guidelines
\end{itemize}
\section{Treatment Options}
\begin{itemize}[noitemsep]
  \item Treatment options include timed bright light therapy in the morning, low-dose melatonin (0.5--3 mg) administered 5--6 hours before target bedtime, and consistent sleep-wake scheduling
  \item Chronotherapy may be considered in refractory cases.
  \item Medications to manage symptoms and address underlying mechanisms
  \item Lifestyle modifications and self-care measures
  \item Physical therapy and rehabilitation
  \item Surgical intervention when indicated
  \item Regular monitoring and follow-up care
\end{itemize}
\section{Living with It}
\begin{itemize}[noitemsep]
  \item Follow your treatment plan as prescribed
  \item Maintain regular follow-up with your healthcare provider
  \item Make necessary lifestyle adjustments
  \item Seek support from family, friends, or support groups
  \item Stay informed about your condition
\end{itemize}
\section{Outlook}
The outlook is generally positive with treatment compliance. The outlook varies depending on the specific condition, severity at diagnosis, and response to treatment. Early intervention generally leads to better outcomes. Many conditions can be effectively managed with appropriate treatment. Regular follow-up care is important for monitoring and treatment adjustment.
